%=====================================================================
%  Curriculum Vitae -- Tamjid Rahat
%  Standalone LaTeX (no third-party class). Compile with pdfLaTeX.
%  Font & layout mirror the moderncv "banking" style
%  (TeX Gyre Pagella serif; Large bold section titles with a rule).
%=====================================================================
\documentclass[11pt,letterpaper]{article}

\usepackage[T1]{fontenc}
\usepackage[utf8]{inputenc}
\usepackage[margin=0.85in]{geometry}

% TeX Gyre Pagella (Palatino-like serif) -- the font used by moderncv banking.
\usepackage{tgpagella}
\usepackage{microtype}

\usepackage{xcolor}
\definecolor{accent}{HTML}{0284C7}   % azure, matches tamjidrahat.github.io
\definecolor{ink}{HTML}{1F2328}
\definecolor{muted}{HTML}{57606A}

\usepackage{enumitem}
\usepackage{titlesec}
\usepackage[colorlinks=true]{hyperref}
\hypersetup{urlcolor=accent, linkcolor=accent}

\color{ink}
\setlength{\parindent}{0pt}
\setlength{\parskip}{0.5em}
\linespread{1.1}

% ---- Section headings: Large bold accent title with a solid rule beneath ----
%      (banking style; small positive gap so the rule never touches the text) ----
\titleformat{\section}
  {\Large\bfseries\color{accent}}
  {}{0em}{}
  [{\vspace{1pt}{\color{accent!70}\titlerule[0.6pt]}}]
\titlespacing*{\section}{0pt}{2.4ex}{1ex}

% ---- Compact bullet lists with an accent marker ----
\setlist[itemize]{leftmargin=1.35em, itemsep=3pt, topsep=3pt, parsep=0pt,
  label=\textcolor{accent}{\footnotesize$\bullet$}}

% ---- Entry header: bold title (left) + muted date (right) ----
\newcommand{\cvhead}[2]{%
  \begin{tabular*}{\linewidth}{@{}l@{\extracolsep{\fill}}r@{}}%
    \textbf{#1} & {\small\color{muted}#2}%
  \end{tabular*}\\[2pt]}

% ---- Publication venue tag & author self-highlight ----
\newcommand{\venue}[1]{\textbf{\textcolor{accent}{#1}}}
\newcommand{\me}{\underline{T. Al Rahat}}

\begin{document}
\pagestyle{empty}

%=====================================================================
%  Header
%=====================================================================
\begin{center}
  {\fontsize{28}{32}\selectfont\bfseries Tamjid Rahat}\\[8pt]
  \normalsize
  \href{mailto:tamjid@ucla.edu}{tamjid@ucla.edu}\quad\textbullet\quad
  \href{https://tamjidrahat.github.io}{tamjidrahat.github.io}\quad\textbullet\quad
  \href{https://www.linkedin.com/in/tamjidrahat}{LinkedIn}\quad\textbullet\quad
  \href{https://scholar.google.com/citations?user=LY1s3bsAAAAJ\&hl=en}{Google Scholar}
\end{center}

%=====================================================================
%  Research Interests
%=====================================================================
\section{Research Interests}
Agentic AI and software security; automated program repair (APR); LLM agents for
vulnerability detection and remediation; agent evaluation (validators, LLM-as-a-judge);
retrieval-augmented generation (RAG) for code; human-in-the-loop learning; formal methods
and program analysis.

%=====================================================================
%  Education
%=====================================================================
\section{Education}
\cvhead{University of California, Los Angeles (UCLA)}{2024}%
Ph.D. in Computer Engineering \quad\textbullet\quad CGPA 4.0/4.0\\[1pt]
\textit{Dissertation:} ``Automated Detection and Mitigation of Vulnerabilities in
Single Sign-on Implementations.''

\cvhead{University of Virginia, Charlottesville, VA}{2020}%
M.S. in Computer Science \quad\textbullet\quad CGPA 3.90/4.0

\cvhead{Bangladesh University of Engineering and Technology (BUET)}{2016}%
B.Sc. in Computer Science and Engineering \quad\textbullet\quad CGPA 3.84/4.0

%=====================================================================
%  Experience
%=====================================================================
\section{Experience}
\cvhead{Amazon Web Services (AWS), San Jose, CA}{06.2024 -- Present}%
\textit{Applied Scientist, AWS Security}
\begin{itemize}
  \item Designed agentic systems for automated program repair (APR) that detect and fix
        security vulnerabilities across heterogeneous services and codebases at production scale.
  \item Shipped production code that improved the precision and recall of agent-generated fixes.
  \item Developed LLM-as-a-judge evaluators to monitor agent output and collect feedback, and
        used that feedback to improve the agents in the next iteration.
  \item Built evaluation benchmarks and regression suites to keep fixes consistent across
        different LLMs and agents.
  \item Optimized the pipeline for scalability and token efficiency, and used
        retrieval-augmented generation (RAG) and coding agents such as Claude and Kiro to
        produce high-quality fixes in production.
\end{itemize}

\cvhead{United International University, Bangladesh}{06.2016 -- 07.2018}%
\textit{Lecturer, Computer Science and Engineering}

\cvhead{Stochastic Logic Limited, Bangladesh}{05.2016 -- 07.2018}%
\textit{Quantitative Software Developer}

%=====================================================================
%  Publications
%=====================================================================
\section{Publications}
\begin{itemize}[itemsep=5pt]
  \item \venue{ASE 2025:} \me{}, Y. Chen, Y. Feng, and Y. Tian. ``Automated Repair of
        OpenID Connect Programs.'' 40th IEEE/ACM Int. Conference on Automated Software
        Engineering.

  \item \venue{ACM CCS 2024:} \me{}, Y. Feng, and Y. Tian. ``AuthSaber: Automated Safety
        Verification of OpenID Connect Programs.'' 31st ACM Conference on Computer and
        Communications Security.

  \item \venue{ACSAC 2023 [Poster]:} M. Xie, \me{}, W. Wang, and Y. Tian. ``Using Program
        Knowledge Graph to Uncover Software Vulnerabilities.'' Annual Computer Security
        Applications Conference.

  \item \venue{ACM CCS 2022:} \me{}, Y. Feng, and Y. Tian. ``Cerberus: Query-driven Scalable
        Vulnerability Detection in OAuth Service Provider Implementations.'' 29th ACM
        Conference on Computer and Communications Security.

  \item \venue{IEEE TPS 2022} [\textbf{Best Paper Award}]: T. Saha, \me{}, N. Aaraj, Y. Tian,
        and N. Jha. ``ML-FEED: Machine Learning Framework for Efficient Exploit Detection.''
        IEEE Int. Conference on Trust, Privacy and Security in Intelligent Systems and
        Applications.

  \item \venue{ACM WPES 2022:} \me{}, M. Long, and Y. Tian. ``Is Your Policy Compliant?
        A Deep Learning-based Empirical Study of Privacy Policies' Compliance with GDPR.''
        21st Workshop on Privacy in the Electronic Society.

  \item \venue{VEHITS 2020:} T. Le, I. ElSayed-Aly, W. Jin, S. Ryu, G. Verrier, \me{},
        B. Park, and Y. Tian. ``Evaluating the Dedicated Short-range Communication for
        Connected Vehicles against Network Security Attacks.'' 6th Int. Conference on
        Vehicle Technology and Intelligent Transport Systems.

  \item \venue{IEEE/ACM ASE 2019:} \me{}, Y. Feng, and Y. Tian. ``OAuthLint: An Empirical
        Study on OAuth Bugs in Android Applications.'' 34th IEEE/ACM Conference on Automated
        Software Engineering.

  \item \venue{ADC 2018:} \me{}, A. Arman, and M. E. Ali. ``Maximizing Reverse k-Nearest
        Neighbors for Trajectories.'' 29th Australasian Database Conference.
\end{itemize}

%=====================================================================
%  Technical Skills
%=====================================================================
\section{Technical Skills}
\begin{itemize}
  \item \textbf{AI / Agentic:} LLMs and frontier models, LLM agents, retrieval-augmented
        generation (RAG), LLM-as-a-judge evaluation, agentic workflow orchestration,
        coding agents (Claude Code, Kiro).
  \item \textbf{Programming Languages:} Python, Java, C/C++, Racket, Intel x86 Assembly, MATLAB.
  \item \textbf{Program Analysis:} WALA, Soot, Ghidra, Pysa, Clang, Z3, Tamarin-prover,
        Rosette, FlowDroid.
  \item \textbf{Machine Learning:} TensorFlow, Keras, PyTorch, scikit-learn.
  \item \textbf{Web \& Data:} JavaScript, PHP, MySQL, SQLite, Oracle.
  \item \textbf{Other:} \LaTeX{}, Bash, HTML, Lisp.
\end{itemize}

%=====================================================================
%  Awards and Grants
%=====================================================================
\section{Awards and Grants}
\begin{itemize}
  \item \textbf{Google Research Paper Award}, 2023.
  \item \textbf{Qualcomm Innovation Fellowship}, Finalist, 2023.
  \item \textbf{Google Vulnerability Reward}, 2022 (\$5{,}000 for reporting security-critical bugs).
  \item \textbf{Vulnerability Research Grant, Google}, 2022 (security research on Google Identity Services).
  \item \textbf{Graduate Education Fellowship}, UCLA, 2022--2023.
  \item \textbf{UVA Computer Science Fellowship}, University of Virginia, 2018--2019.
  \item \textbf{Bug Bounty}, GoFundMe, 2019.
  \item \textbf{Travel Grant}, Amazon Graduate Research Symposium, 2019.
  \item \textbf{Dean's List Award} and \textbf{University Merit Scholarship}, BUET, 2012--2016.
\end{itemize}

%=====================================================================
%  Academic Service
%=====================================================================
\section{Academic Service}
\begin{itemize}
  \item \textbf{Technical Program Committee:} USENIX Security '24 (posters),
        IEEE S\&P '24 (posters), ACM CCS '22 (posters), IEEE S\&P '23 (posters).
  \item \textbf{Session Chair:} ACM CCS '22, ACM WPES '22.
  \item \textbf{Peer Reviewer:} CSF '21, ICSTE '22, JSS '22--23.
  \item \textbf{External Peer Reviewer:} S\&P '21--24, CCS '21--24, USENIX '21--24,
        WWW '22, SenSys '21--22.
\end{itemize}

%=====================================================================
%  Media Coverage
%=====================================================================
\section{Media Coverage}
\begin{itemize}
  \item \textit{High-Severity Bug Reported in Google's OAuth Client Library for Java},
        The Hacker News, May 2022.
\end{itemize}

\end{document}
